% ECONOMICS_PROBLEM: {"id": "D44-83", "title": "Symmetric equilibrium existence with multidimensional auction information", "jel_code": "D44", "source_id": "OP-0592"}
% FIRST_POSTED: 2026-10-09
% ATTEMPT_STATUS: unresolved
% ENTRY_KIND: research_attempt
% !TeX program = pdflatex
\documentclass[11pt,a4paper]{article}
\usepackage[T1]{fontenc}
\usepackage[utf8]{inputenc}
\usepackage{lmodern}
\usepackage[margin=24mm]{geometry}
\usepackage{amsmath,amssymb,amsthm,mathtools}
\usepackage{booktabs,longtable,array,enumitem,xurl,hyperref,listings,microtype}
\hypersetup{colorlinks=true,linkcolor=blue,urlcolor=blue,citecolor=blue,pdftitle={Research proof audit}}
\setlength{\parindent}{0pt}
\setlength{\parskip}{5pt}
\setlength{\emergencystretch}{3em}
\setlist{nosep,leftmargin=2em}
\newtheorem{theorem}{Theorem}[section]
\newtheorem{lemma}[theorem]{Lemma}
\newtheorem{proposition}[theorem]{Proposition}
\newtheorem{corollary}[theorem]{Corollary}
\theoremstyle{remark}\newtheorem{remark}[theorem]{Remark}
\newcommand{\R}{\mathbb R}
\newcommand{\N}{\mathbb N}
\newcommand{\E}{\mathbb E}
\newcommand{\Var}{\operatorname{Var}}
\newcommand{\ind}{\mathbf 1}
\newcommand{\EF}{\mathrm{EF}}
\newcommand{\EFM}{\mathrm{EFM}}
\newcommand{\PO}{\mathrm{PO}}
\newcommand{\MMS}{\mathrm{MMS}}
\DeclareMathOperator*{\argmax}{arg\,max}
\DeclareMathOperator*{\argmin}{arg\,min}
\lstset{basicstyle=\ttfamily\footnotesize,breaklines=true,columns=fullflexible,keepspaces=true,frame=single}
\begin{document}
\begin{center}
{\LARGE\bfseries Symmetric auction equilibrium: exact binary-family analysis}\par\medskip
{\large D44-83.tex}\par
Research attempt and adversarial audit --- 9 October 2026
\end{center}
\label{problem:OP-0592}
\label{supp:auction-existence}
\textbf{Tools.} Web browsing: YES. Computation: YES (Python, exact integer/rational checks, and explicitly identified numerical diagnostics).
\textbf{Status convention.} A full proof of a restricted theorem does not resolve the full input. Literature results are marked as imported; no novelty or exhaustive current-literature certification is claimed. Computational searches certify only their stated finite domains.

\section{FORMAL RESTATEMENT}
An admissible primitive tuple $\mathcal A$ consists of integers $n\ge2$, $1\le k<n$, $S\ge1$, a taste bound $\bar t>0$, strictly positive continuous probability densities $f$ on $[0,1]$ and $w$ on $[0,\bar t]$, and a $C^1$ function $u(q,t)$ strictly increasing in each variable. The common quality $q$ has density $f$. Independently conditional on $q$, bidders draw tastes with density $w$ and signals satisfying
\[
 \Pr(s_i=s\mid q,t_i)=\alpha_s(q)\beta_s(t_i),\qquad
 \sum_{s=1}^S\alpha_s(q)\beta_s(t)=1.
\]
Here $\alpha_s$ is $C^1$ and bounded away from zero, $\beta_s\ge0$ is measurable, and $\alpha_r/\alpha_s$ is strictly increasing in $q$ when $r>s$. The conditional laws are proper, not unnormalized likelihoods. Bidders are unit-demand and risk-neutral: a winner's utility is $u(q,t)$ minus its payment; a loser receives zero. The $k$ highest nonnegative bids win and each pays the $(k+1)$st highest bid; a tie for the remaining units is resolved uniformly among the tied bidders.

For a symmetric profile $b=(b_s)$ define its conditional payoff $\Pi_s(t,x;b)$ at each positive-probability conditional type. A pure Bayesian equilibrium requires, for almost every $(t,s)$ and every $x\ge0$,
\[
 \Pi_s(t,b_s(t);b)\ge\Pi_s(t,x;b).
\]
Write $\mathcal E^{\mathrm c}$ for continuous strictly taste-increasing profiles, $\mathcal E^{\mathrm{A2}}\subseteq\mathcal E^{\mathrm c}$ when each function is differentiable except at finitely many points, and $\mathcal E^{\mathrm{mon}}$ for measurable weakly taste-increasing profiles, allowing jumps and pooling. Null supported types impose no equilibrium restriction. An $\varepsilon$ equilibrium replaces the right side by $\Pi_s(t,x;b)-\varepsilon$.

The four questions are: (a) whether $\forall\mathcal A$, $\mathcal E_{\mathcal A}^{\mathrm{A2}}\ne\varnothing$; (b) classify both continuous and monotone existence for every $T>0$ in the binary family below; (c) whether $\forall\mathcal A$, $\mathcal E_{\mathcal A}^{\mathrm{mon}}\ne\varnothing$; and (d) give primitive restrictions and fixed-$\mathcal A$, every-$\varepsilon>0$ approximation results. A disproof of (a) is not a disproof of (c).

In (b), $n=2$, $k=1$, $q\sim U[0,1]$, $t_i\sim U[0,T]$, $u(q,t)=q+t$, $\beta_1=\beta_2=1$, and
\[
 \alpha_1(q)=\frac34-\frac q2,\qquad \alpha_2(q)=\frac14+\frac q2.
\]
All logarithms are natural. The degenerate $T=0$ is excluded because the stated positive-density taste model has $T>0$. Stress points are global rather than local optimality, possible singular continuous bid functions, common empty bid intervals, and keeping the auction rule unchanged.

\section{QUICK LITERATURE/CONTEXT CHECK}
The input names an 8 October 2026 research note, \emph{Continuous-Strategy Nonexistence in a Two-Bidder Auction}. A separate copy of that note was not available in the eleven attachments, and the search did not independently retrieve it. Its reported threshold is therefore checked below directly from the supplied primitives; no unseen proof is imported.

Jackson's accessible author manuscript \cite{jackson} exhibits nonexistence with discrete tastes, binary quality, and signals that can reveal quality completely. Those features do not satisfy this input's positive continuous densities and everywhere-positive signal likelihoods. It is relevant context, not a counterexample to the present primitive class. No general finite-auction monotone existence theorem is assumed here. The binary-family conclusions below follow from the displayed calculations and global incentive verification.

\section{ATTACK PLAN}
This is an equilibrium-existence and endogenous-information problem. The proof track derives posterior moments, obtains an inverse-bid system without assuming differentiability, and verifies candidate bids by the sign of every payoff increment. The construction track first separates the two bid supports and then removes a middle inverse-curve segment to create a common bid gap. The counterexample track tests the interval between the two reported thresholds and checks whether an argument rules out continuous strategies only, or all monotone strategies.

\begin{center}\begin{tabular}{p{.25\linewidth}p{.66\linewidth}}\toprule
Tool & Role \\\midrule
Bayes' rule & Computes the exact conditional pair-signal values.\\
Single-crossing payoff increments & Turns a marginal identity into a global best response.\\
Stieltjes measures & Handles continuous strictly increasing bids even when derivatives fail.\\
Continuous bounded-variation calculus & Justifies the inverse-curve differential identities.\\
Separation of bid supports & Solves the small-taste range explicitly.\\
Reciprocal symmetry & Matches the two taste cutoffs at a common discontinuity.\\
Symbolic algebra and quadrature & Checks constants and the critical integral.\\
Coupling and uniform continuity & Smooths jump strategies at fixed $n$ for $\varepsilon$ equilibria.\\\bottomrule
\end{tabular}\end{center}
The selected path completely treats the binary family, disproves universal A2 existence, and gives two restricted answers to (d). General question (c) is kept separate.

\section{WORK}
\subsection{Posterior moments and a payoff formula}
Set
\[
 a=\frac9{26},\qquad d=\frac{17}{26},\qquad \delta=\frac2{13},
 \qquad a=\frac12-\delta,\quad d=\frac12+\delta.
\]
Direct polynomial integration gives
\[
 \left(\int_0^1\alpha_s(q)\alpha_r(q)\,dq\right)_{s,r}
 =\frac1{48}\begin{pmatrix}13&11\\11&13\end{pmatrix},\qquad
 \left(\int_0^1q\alpha_s(q)\alpha_r(q)\,dq\right)_{s,r}
 =\frac1{96}\begin{pmatrix}9&11\\11&17\end{pmatrix}.
\]
Each own signal has probability $1/2$. Consequently, with
\[
 (c_{sr})=\begin{pmatrix}13&11\\11&13\end{pmatrix},\qquad
 (\mu_{sr})=\begin{pmatrix}a&1/2\\1/2&d\end{pmatrix},
\]
we have $\Pr(s_{-i}=r\mid s_i=s)=c_{sr}/24$ and
$\E[q\mid s_i=s,s_{-i}=r]=\mu_{sr}$. The opponent's taste remains uniform and independent of these signals.

Against strictly increasing strategies, including strategies with jumps but no flat taste intervals, the opponent's bid distribution has no atoms. For every $t$ and $x\ge0$,
\begin{equation}\label{eq:auctionpay}
 \Pi_s(t,x;b)=\frac1{24T}\sum_{r=1}^2 c_{sr}
   \int_{\{z:b_r(z)<x\}}\bigl(t+\mu_{sr}-b_r(z)\bigr)\,dz.
\end{equation}
This follows by conditioning on the opponent's signal and taste; the second price is exactly its bid. Ties have probability zero. The formula is continuous in $(t,x)$ when $b_r$ is strictly increasing, since the induced bid laws have no atoms and all their bids are bounded on the compact taste interval. Therefore, for continuous candidate strategies, an almost-everywhere best-response inequality extends to every type: first use a countable dense set of deviations $x$, extend in $t$ by continuity, and then extend in $x$.

\begin{proposition}[Separated-support equilibria]\label{prop:separated}
If $0<T\le\delta$, then
\[
 b_1(t)=t+a,\qquad b_2(t)=t+d
\]
is a symmetric equilibrium in $\mathcal E_T^{\mathrm{A2}}$.
\end{proposition}
\begin{proof}
The low-signal support ends at $T+a\le1/2$, below the high-signal support beginning at $d$. Against its own signal, either type's surplus from an opponent of taste $z$ changes sign at $z=t$. A low-signal type's surplus from a high-signal opponent is
$t+1/2-(z+d)=t-\delta-z\le0$; it optimally wins none of that support. A high-signal type's surplus from a low-signal opponent is
$t+1/2-(z+a)=t+\delta-z\ge0$; it optimally wins all of that support. The proposed bids implement both desired threshold decisions simultaneously. Integrating these pointwise surplus comparisons in \eqref{eq:auctionpay} verifies every bid deviation, including deviations into the support gap or beyond both supports.
\end{proof}

\subsection{A derivative-free necessary condition for continuous equilibrium}
\begin{lemma}[Inverse-bid identities]\label{lem:inverse}
Suppose $b_1,b_2$ are continuous strictly increasing equilibrium strategies and their bid supports overlap on an open interval $I$. Let $t_r(b)$ be their continuous increasing inverses on $I$ and let $dt_r$ denote their Stieltjes measures. Then
\begin{align}
 13(t_1+a-b)\,dt_1+11(t_1+1/2-b)\,dt_2&=0,\label{eq:stieltjes1}\\
 11(t_2+1/2-b)\,dt_1+13(t_2+d-b)\,dt_2&=0.\label{eq:stieltjes2}
\end{align}
In particular $t_1(b)-t_2(b)\ge\delta$ on $I$.
\end{lemma}
\begin{proof}
Fix $x<y$ in a compact subinterval of $I$. Optimality at $x$ of the type $t_s(x)$ gives
\[
 I_x:=\sum_r c_{sr}\int_{(x,y)}(t_s(x)+\mu_{sr}-b)\,dt_r(b)\le0.
\]
Optimality at $y$ of the type $t_s(y)$ gives the same integral with $t_s(y)$ in place of $t_s(x)$, denoted $I_y$, and $I_y\ge0$. Let $M_s=\sum_r c_{sr}dt_r$. The difference is
$I_y-I_x=(t_s(y)-t_s(x))M_s((x,y))$. The integral $J_{x,y}$ with $t_s(b)$ inside the integrand lies between $I_x$ and $I_y$. Hence
\[
 |J_{x,y}|\le(t_s(y)-t_s(x))M_s((x,y)).
\]
Partition any fixed compact interval into increasingly fine bid intervals. Since $t_s$ is uniformly continuous there and $M_s$ is finite, summing this bound shows that the integral of the signed measure defining $J$ is zero on every compact interval. Endpoints carry no Stieltjes mass. This proves \eqref{eq:stieltjes1}--\eqref{eq:stieltjes2} as identities of signed measures; neither inverse is assumed absolutely continuous.

Put $\nu=dt_1+dt_2$ and $z_1=1/2+t_1-b$. The first identity is
\[
 13(z_1-\delta)dt_1+11z_1dt_2=0.
\]
Both Stieltjes measures are positive and $\nu$ assigns positive mass to every open subinterval. If the continuous function $z_1$ were negative, both coefficients would be strictly negative in a neighborhood, contradicting this identity. If $z_1>\delta$ the analogous positive-coefficient contradiction applies. Thus $0\le z_1\le\delta$. Rearranging the identity now gives
\[
 dt_1=p\,d\nu,\qquad dt_2=(1-p)\,d\nu,
 \qquad p=\frac{11z_1}{2(1-z_1)}\in[0,1].
\]
Here $p$ is continuous. Substitution in both identities yields
\[
 z_1=\frac{2p}{11+2p},\qquad
 z_2:=b-t_2-\frac12=\frac{2(1-p)}{13-2p}.
\]
The equalities initially hold $\nu$-almost everywhere; continuity and full support extend them everywhere. Finally
\[
 t_1-t_2-\delta=z_1+z_2-\delta
   =\frac{96p(1-p)}{13(13-2p)(11+2p)}\ge0.
\]
\end{proof}

\begin{lemma}[Support order and the only possible overlap length]\label{lem:length}
For a continuous strictly increasing equilibrium, write $L_s=b_s(0)$ and $U_s=b_s(T)$. Then $L_1<L_2$ and $U_1<U_2$. Either the supports are disjoint (possibly touching) and $T\le\delta$, or they overlap and
\[
 T\ge T_c:=\frac16+\frac{11}{144}\log\frac{13}{11}.
\]
\end{lemma}
\begin{proof}
On an open bid interval containing only signal $s$'s bids, local deviations between $b_s(z)$ are enough to show $b_s(t)=t+\mu_{ss}$ for every type bidding in its interior. Indeed its variable payoff is a positive constant times
$\int^z(t+\mu_{ss}-b_s(v))\,dv$, a differentiable function of $z$; at the interior optimum $z=t$ its derivative is zero. This argument uses continuity of $b_s$, not differentiability of bids.

If $L_1>L_2$, the high-signal-only bottom interval gives $L_2=d$. A low-signal type of taste zero bidding $L_1$ wins a positive mass of high-signal opponents with prices at least $d>1/2$, and wins no positive mass of its own signal. Its payoff is strictly negative; bidding zero loses surely and gives zero. This is a contradiction. If $L_1=L_2$, an overlap immediately above that endpoint has $t_1,t_2\to0$, contradicting Lemma~\ref{lem:inverse}. Therefore $L_1<L_2$.

If $U_1>U_2$, a low-signal-only top interval gives $U_1=T+a$. A high-signal type of taste $T$ can profit by raising its bid above $U_1$: every newly won low-signal opponent in that top interval has surplus
$T+1/2-(z+a)=T-z+\delta>0$. It gives up no winnings and pays the opponent's bid. Equal upper endpoints contradict Lemma~\ref{lem:inverse} as $t_1,t_2\to T$. Thus $U_1<U_2$.

If $L_2\ge U_1$, the single-signal argument throughout both interiors gives $b_1(t)=t+a$ and $b_2(t)=t+d$. When $T>\delta$, a low-signal type at $T$ can raise its bid to win a positive sufficiently small interval of high-signal tastes $z<T-\delta$, all of positive surplus $T-\delta-z$. It was already winning all its own signal. Thus disjoint supports require $T\le\delta$.

It remains to treat the overlap $I=(L_2,U_1)$. In Lemma~\ref{lem:inverse}, $p$ is locally of bounded variation: $z_1=1/2+t_1-b$ has that property, and the smooth map $11z/[2(1-z)]$ has bounded derivative on $[0,\delta]$. Where $p>0$ put $r=(1-p)/p$. The inverse identities become
\begin{equation}\label{eq:inverseparam}
 t_1=b-\frac12+\frac{2}{13+11r},\qquad
 t_2=b-\frac12-\frac{2r}{11+13r},\qquad dt_2=r\,dt_1.
\end{equation}
Their continuous bounded-variation chain rule gives
\[
 (1-r)\bigl(db-h(r)\,dr\bigr)=0,
 \qquad
 h(r)=\frac{22(169r^2+334r+169)}{(13+11r)^2(11+13r)^2}>0.\tag{*}
\]
Here and below $db$ is Lebesgue measure on the bid axis. To justify the chain rule being used, for a continuous bounded-variation function $r$ on a compact interval and a $C^1$ function $F$, telescope $F(r)$ over partitions. The error between each increment and $F'(r)$ times the increment of $r$ is bounded by the oscillation of $F'$ over that increment times its absolute variation. Uniform continuity of $r$ and $F'$ makes the total error at most a vanishing quantity times the finite total variation. This proves $d(F(r))=F'(r)dr$ as measures. Applying it to the two rational functions in \eqref{eq:inverseparam}, and clearing denominators, gives $(*)$.

Let $H(r)=\int_0^r h(v)\,dv$. On any interval with finite $r\ne1$, $(*)$ implies $b-H(r)$ is constant, so $r$ is strictly increasing there. A point with $r=0$ cannot lie in the interior of $I$: continuity gives a neighborhood with $r<1$ and the same identity, forcing negative $r$ immediately to its left. A point with $p=0$, corresponding to $r=\infty$, is likewise impossible. To see this without manipulating infinity, put $\rho=p/(1-p)$ near such a point. Directly using \eqref{eq:inverseparam} with $\rho=1/r$ gives
\[
 (1-\rho)\bigl(db+h(\rho)\,d\rho\bigr)=0.
\]
Thus $b+H(\rho)$ is constant near $\rho=0$, forcing negative $\rho$ to its right. Consequently $0<r<\infty$ everywhere in $I$.

The single-signal bottom segment gives $b_1(t)=t+a$ up to $L_2$. At the lower overlap endpoint $t_2\to0$ and $1/2+t_1-b\to\delta$, so \eqref{eq:inverseparam} gives
\[
 r\to0,\quad L_2=1/2,\quad t_1(L_2)=\delta.
\]
The single-signal top segment gives $b_2(t)=t+d$ from $U_1$ upward. At the upper overlap endpoint, the identities similarly give
\[
 r\to\infty,\quad U_1=T+1/2,\quad t_2(U_1)=T-\delta.
\]
The continuous function $r$ must therefore pass from $0$ through $1$ to $\infty$. Its level set $r=1$ is a closed interval, possibly a single point: otherwise a component of its complement between two level-one points would be strictly increasing yet have equal endpoint limits. Before and after that level interval, $b-H(r)$ is constant. Denote the length of the level interval in bid coordinates by $\ell\ge0$. The total overlap width is therefore
\[
 U_1-L_2=H(\infty)+\ell.
\]
The width on the left is $T$. An explicit antiderivative, normalized by $H(0)=0$, is
\begin{equation}\label{eq:H}
 H(r)=-\frac{145r+143}{858r^2+1740r+858}
 +\frac{11}{288}\log\!\frac{r+11/13}{r+13/11}+T_c.
\end{equation}
Differentiating the displayed rational and logarithmic terms gives $h$; substitution at zero gives zero, and the limit at infinity is $T_c$. Hence $T=T_c+\ell\ge T_c$.
\end{proof}

\subsection{Construction and global verification for every taste range}
Define, for $r\in[0,\infty)$,
\[
 B(r)=\frac12+H(r),\qquad
 X(r)=H(r)+\frac2{13+11r},\qquad
 Y(r)=H(r)-\frac{2r}{11+13r}.
\]
Differentiating and putting terms over a common denominator gives
\begin{equation}\label{eq:XYderivatives}
 X'(r)=\frac{1056(r+1)}{(13+11r)^2(11+13r)^2}>0,
 \qquad Y'(r)=rX'(r).
\end{equation}
Thus $X$ and $Y$ are strictly increasing on their ranges (the single point $Y'(0)=0$ does not create a constant interval). Their endpoint values are
\[
 (X(0),Y(0))=(\delta,0),\qquad
 (X(\infty),Y(\infty))=(T_c,T_c-\delta).
\]
The parameterized curve $(b,t_1,t_2)=(B(r),X(r),Y(r))$ satisfies both inverse-bid identities.

\begin{proposition}[Continuous existence for $T\ge T_c$]\label{prop:continuouslarge}
For every $T\ge T_c$ there is an equilibrium in $\mathcal E_T^{\mathrm{A2}}$.
\end{proposition}
\begin{proof}
For $T=T_c$ use the curve just defined on the overlapping bid interval. Append the low-signal-only segment $b_1(t)=t+a$ for $0\le t\le\delta$ and the high-signal-only segment $b_2(t)=t+d$ for $T-\delta\le t\le T$. The endpoint values match: the lower common bid is $1/2$, the upper common bid is $T+1/2$.

For $T>T_c$ put $\ell=T-T_c$. Follow the same curve until $r=1$; insert an interval of bid length $\ell$ on which
\[
 t_1=b-5/12,\qquad t_2=b-7/12;
\]
then continue the $r\ge1$ branch after adding $\ell$ to each of $b,t_1,t_2$. At $r=1$ these formulas coincide with the endpoints of the original curve. On the inserted interval $dt_1=dt_2=db$, and substituting into \eqref{eq:stieltjes1}--\eqref{eq:stieltjes2} gives zero. Append the same two outer segments with the new value of $T$.

Both inverse functions increase strictly across their full taste ranges. Their inverses are continuous and strictly increasing. On the interiors of the finitely many segments, the derivatives in \eqref{eq:XYderivatives} are positive wherever an inverse derivative is needed; the only possible exceptions occur at their finitely many junctions and endpoints. Thus the resulting bids are in the A2 regularity class. Global optimality is proved in Proposition~\ref{prop:global}; this construction satisfies all of that proposition's hypotheses.
\end{proof}

\begin{proposition}[A common bid gap for $\delta<T<T_c$]\label{prop:jump}
For every $\delta<T<T_c$, there is a symmetric strictly taste-increasing measurable equilibrium with one jump in each signal's bid function and no bid atoms.
\end{proposition}
\begin{proof}
For $z\in(0,1)$ set
\[
 D(z)=X(1/z)-X(z),\qquad T(z)=T_c-D(z).
\]
Because $X$ is strictly increasing and continuous, $D$ is strictly decreasing and continuous in $z$. Its endpoint limits are $T_c-\delta$ as $z\downarrow0$ and zero as $z\uparrow1$. Thus $T(z)$ runs strictly increasingly and continuously through $(\delta,T_c)$. Fix the unique $z$ with $T(z)=T$.

The taste difference
\[
 X(r)-Y(r)=\frac2{13+11r}+\frac{2r}{11+13r}
\]
is unchanged when $r$ is replaced by $1/r$ (cross-multiplication verifies this identity). It follows that
\[
 Y(1/z)-Y(z)=D(z).
\]
Keep the branch $0\le r\le z$ of $(B,X,Y)$ unchanged. Delete the branch $z<r<1/z$. Keep $r\ge1/z$ after subtracting $D(z)$ from all three coordinates. At the two retained endpoints, the taste coordinate for each signal is the same, but the bid coordinate jumps by
\[
 G(z)=B(1/z)-D(z)-B(z)
   =\frac2{13+11z}-\frac2{13+11/z}>0.
\]
This defines a common empty bid interval. The low-signal jump occurs at taste $X(z)$; the high-signal jump occurs at taste $Y(z)$. Each cutoff is strictly inside its relevant overlap taste range: for example $T-X(z)=T_c-X(1/z)>0$, and the high-signal inequality follows from the corresponding identity for $Y$.

Append $b_1(t)=t+a$ below taste $\delta$ and $b_2(t)=t+d$ above taste $T-\delta$. All endpoints fit because subtracting $D$ changes the upper inverse limits to $T$ and $T-\delta$. Invert each strictly increasing retained branch. At a jump cutoff assign either endpoint bid, keeping monotonicity. A single cutoff taste has probability zero, so no bid atom is created. Shifting $b,t_1,t_2$ by the same number preserves every expression $t_s+\mu_{sr}-b$, and hence preserves the inverse-bid identities. Proposition~\ref{prop:global} proves optimality, including at the two cutoff tastes.
\end{proof}

\begin{proposition}[Global, not merely local, incentives]\label{prop:global}
The profiles constructed in Propositions~\ref{prop:continuouslarge} and \ref{prop:jump} are Bayesian equilibria against every nonnegative bid deviation and for every taste of either signal.
\end{proposition}
\begin{proof}
On every active overlapping bid segment, let $t_r(b)$ be the inverse taste coordinate. In the jump construction, extend both inverse functions as constants across the common bid gap. The gap has no opponent bid mass. Write $M_s=\sum_r c_{sr}dt_r$, a nonnegative measure. The inverse identities, subtracted from the payoff increment for an arbitrary own taste $t$, give
\begin{equation}\label{eq:globalincrement}
 d\Pi_s(t,b)=\frac{t-t_s(b)}{24T}\,dM_s(b)
\end{equation}
on the overlapping segments. The inverse coordinate $t_s(b)$ is increasing, with only the declared constant gap. Consequently these increments are nonnegative before the bid or gap corresponding to taste $t$, and nonpositive afterwards. The payoff is constant on the common gap. A cutoff taste is therefore indifferent among every bid in that gap, including both endpoint choices.

On the low-signal-only bottom segment the opponent bid is $z+a$ for $0\le z\le\delta$. A low-signal type's increment has the sign of $t-z$. A high-signal type's increment has the sign of $t+\delta-z$, which is nonnegative for every taste $t\ge0$ on that entire segment. Thus each high-signal type prefers winning all of that bottom segment; a low-signal type whose own bid lies there has the exact threshold $z=t$.

On the high-signal-only top segment the opponent bid is $z+d$ for $T-\delta\le z\le T$. A high-signal type's increment has the sign of $t-z$. A low-signal type's increment has the sign of $t-\delta-z\le0$. Thus every low-signal type prefers winning none of the top segment, and high-signal types with bids there again choose the exact taste threshold. These signs also show that types outside an overlap taste range never improve by entering it: in \eqref{eq:globalincrement} their taste lies entirely below or entirely above its inverse coordinates.

Below the smallest opponent bid or above the largest, moving one's bid changes no opponent mass; the payoff is constant there. Combining the signs on all consecutive segments proves that the submitted bid is a global maximizer of \eqref{eq:auctionpay}. This checks arbitrary deviations, not just adjacent or stationary ones. All constructed distributions are atomless, so the specified tie rule agrees with this calculation.
\end{proof}

\begin{theorem}[Complete classification of the binary family]\label{thm:classification}
For every $T>0$,
\[
 \mathcal E_T^{\mathrm c}\ne\varnothing
 \quad\Longleftrightarrow\quad
 \mathcal E_T^{\mathrm{A2}}\ne\varnothing
 \quad\Longleftrightarrow\quad
 T\in(0,2/13]\ \cup\ [T_c,\infty).
\]
Moreover $\mathcal E_T^{\mathrm{mon}}\ne\varnothing$ for every $T>0$.
\end{theorem}
\begin{proof}
The two possible continuous-support configurations and Lemma~\ref{lem:length} prove necessity even without differentiability. Proposition~\ref{prop:separated} supplies the small range, and Proposition~\ref{prop:continuouslarge}, verified by Proposition~\ref{prop:global}, supplies the large range. In the intermediate range, Proposition~\ref{prop:jump} and the same global verification supply a monotone equilibrium. The A2 class is a subset of the continuous class, while the constructed continuous equilibria are A2, proving all asserted equivalences.
\end{proof}

\begin{corollary}[Explicit disproof of question (a)]\label{cor:auctioncounter}
The binary primitives above with $T=1/6$ satisfy all the stated assumptions but have $\mathcal E_T^{\mathrm{A2}}=\varnothing$; in fact $\mathcal E_T^{\mathrm c}=\varnothing$.
\end{corollary}
\begin{proof}
The uniform densities are strictly positive and continuous on their supports. The two likelihoods are $C^1$, take values in $[1/4,3/4]$, and sum to one. Their ratio satisfies
\[
 \frac{d}{dq}\frac{\alpha_2(q)}{\alpha_1(q)}
  =\frac{1}{2\alpha_1(q)^2}>0.
\]
The factors $\beta_s=1$ are admissible. Signals can be generated independently conditional on $q$, and tastes independently of $q$ and of all such signal draws. The value $q+t$ is $C^1$ and strictly increasing in each coordinate. With $n=2,k=1$, the stated payment and tie rules are exactly those used above. Finally
\[
 \frac2{13}<\frac16<\frac16+\frac{11}{144}\log\frac{13}{11},
\]
since $12<13$ and $\log(13/11)>0$. The classification gives nonexistence in the asserted regularity class. This counterexample does not disprove question (c); its monotone equilibrium was constructed explicitly.
\end{proof}

\subsection{Two fully proved, but restricted, answers to variant (d)}
\begin{proposition}[Fixed-binary-auction continuous approximation]
For each fixed $T>0$ and each $\varepsilon>0$, the binary family has a symmetric strictly taste-increasing continuous A2 profile that is an $\varepsilon$ Bayesian equilibrium for every type and every nonnegative bid deviation.
\end{proposition}
\begin{proof}
When the exact continuous equilibria exist, use them. Otherwise take the one-jump profile of Proposition~\ref{prop:jump}. For each signal $s$, write $c_s$ for its cutoff taste and $[L,U]$ for the common empty bid interval. Choose $\eta>0$ smaller than the distance of either cutoff to the taste endpoints and to its neighboring segment junctions. Replace $b_s$ on $[c_s-\eta,c_s+\eta]$ by the straight line joining its values at those two endpoints; keep it unchanged elsewhere. The resulting $b_s^\eta$ is continuous and strictly increasing, and piecewise differentiable with finitely many exceptions.

All old and new bids lie in a fixed bounded interval $[0,B_*]$, independent of sufficiently small $\eta$. Conditional on either own signal, couple the old and new opponents using the same quality, taste and signal. Their bids differ only when that opponent's taste belongs to the length-$2\eta$ band for its signal, an event of conditional probability at most $2\eta/T$. For any own $t,x$, a realized payoff has absolute value at most $K=1+T+B_*$. Thus the uniform payoff difference is bounded by
\[
 \sup_{s,t,x\ge0}|\Pi_s(t,x;b^\eta)-\Pi_s(t,x;b)|
 \le\frac{4K\eta}{T}.\tag{**}
\]
The bound covers arbitrarily large deviations because the own bid is never the second price when it wins.

Let $V_s(t)=\sup_{x\ge0}\Pi_s(t,x;b)$. It is enough to maximize over $[0,B_*+1]$, since all larger bids win the same opponent mass. Each payoff is continuous in $(t,x)$ and is $1$-Lipschitz in $t$; hence $V_s$ is also $1$-Lipschitz. Outside the interpolation band the old chosen bid is exactly optimal against $b$. Inside the band, its regret against $b$ converges uniformly to zero. To prove that last assertion, otherwise choose a sequence $\eta_j\downarrow0$, tastes $t_j$ in the band, and regrets bounded below by some positive constant. A subsequence of the interpolated bids converges to a point $x\in[L,U]$, while $t_j\to c_s$. Every such $x$ is optimal for the cutoff type by Proposition~\ref{prop:global}. Continuity of $\Pi_s$ and $V_s$ then makes those regrets tend to zero, a contradiction. There are only two signals, so convergence is uniform over both.

Finally the regret against $b^\eta$ is at most the regret against $b$ plus twice the uniform error in $(**)$. Taking $\eta$ sufficiently small makes their sum at most $\varepsilon$. Both the primitives and the finite number of bidders stay fixed throughout.
\end{proof}

\begin{proposition}[An uninformative-signal sufficient condition]
For any fixed $n,k$, suppose $S=1$ and
\[
 v(t):=\int_0^1u(q,t)f(q)\,dq\ge0\quad(0\le t\le\bar t).
\]
Then $b(t)=v(t)$ is a symmetric A2 equilibrium under the input's auction rule.
\end{proposition}
\begin{proof}
With one signal, the normalization makes the observed signal uninformative. Opponents' tastes, and hence bids depending on their tastes, are independent of $q$. Given any opponents' pure taste strategies, let $Z$ be their $k$th highest bid. Winning above $Z$ makes the payment $Z$, independent of one's submitted winning bid. Conditional on that threshold and own taste, the expected value is $v(t)$, so the incremental payoff from winning is $v(t)-Z$. Bidding $v(t)$ wins exactly when this is positive, loses when it is negative, and is indifferent at equality. Thus it is a best response to any such opponents' strategies; symmetric play is an equilibrium. The integral is $C^1$ by continuity of the derivatives on the compact domain. It is strictly increasing because $u(q,t')-u(q,t)>0$ for every $q$ whenever $t'>t$, and $f$ is positive. The added nonnegativity condition makes the bid admissible. This proves the claimed A2 property as well as global incentives.
\end{proof}

\section{VERIFICATION}
\textbf{Exact algebra.} The accompanying Python program checks the pair-signal integrals, the identities $H'=h$, \eqref{eq:XYderivatives}, and the reciprocal taste-gap relation using symbolic arithmetic. They can also be verified directly by clearing the displayed positive denominators. The signs $X'>0$, $Y'>0$ for $r>0$, and $h>0$ control every inverse and uniqueness claim.

\textbf{Numerical diagnostic, not a certificate of exact equilibrium.} The program obtains
\[
 T_c\approx0.1794277425784363,\qquad
 \int_0^\infty h(r)\,dr\approx0.17942774257843636.
\]
For $T=1/6$, solving $T(z)=T$ gives approximately
\[
 z=0.3328174256,\quad D=0.0127610759,\quad G=0.0766109090.
\]
The cutoff tastes are approximately $0.1650687527$ for signal 1 and $0.0015979140$ for signal 2. A discretized diagnostic using 10,000 opponent tastes per signal and 61 grid types plus two types immediately adjacent to each signal's cutoff found a maximum positive computed deviation residual of about $4.09\cdot10^{-6}$. This nonzero finite-grid residual is not rounded into an exact-zero claim; the exact certificate is Proposition~\ref{prop:global}. Its bound does not rely on quadrature or on an inability to find a profitable grid deviation.
\begin{lstlisting}
compute posterior pair-signal probabilities and means exactly
verify H' = h, X' > 0, and Y' = r*X'
if T <= 2/13:
    use the two separated affine strategies
elif T >= Tc:
    use the inverse curve, inserting length T-Tc at r=1
else:
    solve Tc - (X(1/z)-X(z)) = T for z in (0,1)
    delete the middle branch; shift the upper branch by -D
for each signal and tested taste:
    compare the prescribed payoff to all grid bid thresholds
# Exact global incentives are proved by signed payoff increments.
\end{lstlisting}

\textbf{Adversarial checks.} The continuous nonexistence proof does not divide by a potentially zero ordinary bid derivative: its inverse equations are Stieltjes identities, and the $r=1$ exceptional set is retained as a possible interval rather than discarded. The possibilities $r=0,\infty$ in the overlap interior are separately excluded. Endpoint best responses follow from continuity and full taste support. The discontinuous profile has no pooling and no bid atoms; assigning a cutoff type either gap endpoint is valid because the entire common gap gives the same payoff. The weak and strict inequalities at $T=\delta$ and $T=T_c$ are checked by the explicit constructions. The approximation proof varies the strategy only, not the primitives, the number of bidders, or the tie rule. The one-signal sufficient condition explicitly adds the nonnegative expected-value hypothesis missing from the general primitive list.

\section{FINAL}
\textbf{Overall LABEL: UNRESOLVED.}

\textbf{Question (a): LABEL: FULL SOLUTION. SUBLABEL: COUNTEREXAMPLE/DISPROOF.}
The exact primitives in Corollary~\ref{cor:auctioncounter}, with $T=1/6$, disprove universal existence in A2, and even in the larger continuous strictly increasing class.

\textbf{Question (b): LABEL: FULL SOLUTION. SUBLABEL: FULL PROOF.}
Theorem~\ref{thm:classification} proves the reported continuous threshold classification without assuming differentiable strategies and proves monotone pure existence for every $T>0$. Both the continuous and discontinuous constructions satisfy global Bayesian incentive inequalities.

\textbf{Question (d), restricted subclaims: FULL PROOF.}
Every fixed member of the binary family has continuous A2 $\varepsilon$ equilibria for every $\varepsilon>0$. Uninformative one-signal primitives with nonnegative expected values have an exact A2 equilibrium for every allowed $n,k$. These results do not characterize all primitives.

\textbf{Exact first gap.} For an arbitrary admissible informative primitive tuple, no argument here proves or disproves the existence of a symmetric measurable weakly taste-increasing pure equilibrium under the unchanged uniform-price and symmetric tie rules. The general fixed-auction approximation and primitive-characterization variants are also not completed.

\textbf{Three next moves.} (1) Extend the sign-of-increments construction to more than two signal classes, allowing endogenous common gaps and possible atoms. (2) Prove a monotone-strategy compactness result that controls tie payoffs at atoms instead of changing the tie rule. (3) Search three-signal or multi-unit primitives for an incompatibility among inverse-bid branches and check all nonlocal deviations, not only first-order conditions.

\textbf{What a counterexample to (c) would need.} It cannot be any member of the fully solved binary family, regardless of its positive taste bound. It must use genuinely different informative primitives, a different number of signals, or a different $n,k$, and must rule out jump and pooling equilibria rather than just continuous profiles. No claim that three signals or more bidders are necessary is made: another two-signal family might behave differently.

\section{Completion Estimate}
60\%

The continuous-existence subquestion is disproved and the entire specified binary family, including monotone existence and fixed-auction continuous approximation, is treated. The general monotone pure-existence question and a general primitive characterization remain unresolved.

This is a subjective estimate of the fraction of the \emph{whole requested mathematical scope} addressed by this report, including the named variants. It is not a probability of correctness, an objectively measured fraction of a proof, or an estimate that the remaining work takes the complementary fraction of the time. Only the eleven supplied files are assessed; no missing rank-1--200 statements are inferred.

\begin{thebibliography}{99}
\bibitem{jackson} M. O. Jackson. \emph{Non-Existence of Equilibrium in Vickrey, Second-Price, and English Auctions}. Review of Economic Design 13 (2009), 137--145. Author manuscript dated 21 September 2005, especially Section 2. \url{https://web.stanford.edu/~jacksonm/nonexist.pdf}.
\bibitem{inputnote} \emph{Continuous-Strategy Nonexistence in a Two-Bidder Auction}, 8 October 2026. Source note named in the supplied problem statement; a separate copy was not independently retrieved. The proof in this report uses the supplied explicit primitives, not an unseen note.
\end{thebibliography}
\end{document}
